\chapter{Spacetime Causal Structure}
\label{ch:causal-structure}

We work in units $\hbar = c = G = 1$ and use the mostly plus metric
signature $g_{\mu\nu} = (-,+,+,+)$ throughout this part.

\section{The Schwarzschild Metric and the Tortoise Coordinate}

The Schwarzschild metric outside a spherically symmetric mass
distribution is
\begin{equation}
	ds^2 = -\left(1-\frac{2M}{r}\right)dt^2
	+ \frac{dr^2}{1-\frac{2M}{r}}
	+ r^2\,d\Omega^2, \qquad r > 2M.
\end{equation}
There is an apparent singularity at $r = 2M$, but this is a coordinate
artifact rather than a true curvature singularity. To see this, we
introduce the tortoise coordinate $r_*$, defined by
\begin{equation}
	r_* = \int\frac{dr}{1-\frac{2M}{r}}
	= r + 2M\ln\!\left(\frac{r}{2M}-1\right),
	\qquad 2M < r < \infty,\quad
	-\infty < r_* < +\infty.
\end{equation}
The name comes from the fact that an infinite coordinate interval
$r_*\to-\infty$ corresponds to approaching the horizon $r\to 2M$,
so the horizon recedes to minus infinity in this coordinate. In terms
of $r_*$, the $t$-$r$ part of the metric becomes
\begin{equation}
	-\left(1-\frac{2M}{r}\right)dt^2
	+ \frac{dr^2}{1-\frac{2M}{r}}
	= \left(1-\frac{2M}{r}\right)(-dt^2 + dr_*^2),
\end{equation}
which is conformally flat, with an overall factor that vanishes at the
horizon but leaves the conformal structure intact. Introducing null
coordinates
\begin{equation}
	u = t - r_*, \qquad v = t + r_*,
\end{equation}
the metric takes the form
\begin{equation}
	ds^2 = -\left(1-\frac{2M}{r}\right)du\,dv + r^2\,d\Omega^2.
\end{equation}

\section{Kruskal Coordinates and the Maximal Extension}

The singularity at $r=2M$ in the original coordinates is removed by
a further change of coordinates. Any transformation of the form
$\tilde u = f(u)$, $\tilde v = g(v)$ preserves the null structure of
the metric. We choose
\begin{equation}
	\tilde u = -4M e^{-u/4M}, \qquad
	\tilde v = 4M e^{v/4M},
\end{equation}
which maps $u\in(-\infty,+\infty)$ to $\tilde u\in(-\infty,0)$ and
$v\in(-\infty,+\infty)$ to $\tilde v\in(0,+\infty)$. In these Kruskal
coordinates the metric becomes
\begin{equation}
	ds^2 = -\frac{2M}{r(\tilde u,\tilde v)}
	e^{-r(\tilde u,\tilde v)/2M}\,d\tilde u\,d\tilde v
	+ r^2(\tilde u,\tilde v)\,d\Omega^2,
\end{equation}
where $r(\tilde u,\tilde v)$ is determined implicitly by
\begin{equation}
	\tilde u\,\tilde v = 4M^2\!\left(1-\frac{r}{2M}\right)
	e^{-r/2M}.
\end{equation}
The metric is now manifestly regular at $r=2M$: when $r=2M$ either
$\tilde u$ or $\tilde v$ vanishes, but the prefactor
$e^{-r/2M}$ is finite. The true singularity at $r=0$ corresponds to
$\tilde u\,\tilde v = 4M^2$, a hyperbola in the $(\tilde u,\tilde v)$
plane. The singularity is genuine because the curvature invariants
diverge there, and geodesics reach it in finite affine parameter.

% [FIGURE: Kruskal diagram showing the four regions, the horizon at
%  tilde{u}=0 or tilde{v}=0, and the singularity at tilde{u}tilde{v}=4M^2]

The Kruskal extension reveals four distinct regions of the maximally
extended Schwarzschild spacetime. Region I is the exterior of the black
hole, $r > 2M$, where $t = \text{const}$ lines converge at spacelike
infinity and $r = \text{const}$ lines are timelike surfaces. Region II
is the black hole interior: once inside, the roles of $r$ and $t$ are
exchanged, with $r = \text{const}$ becoming spacelike surfaces and an
observer being unable to remain at constant $r$. Regions III and IV are
the white hole interior and the second exterior region respectively;
they are physically irrelevant for realistic gravitational collapse but
are part of the maximal extension.

\section{The Penrose Diagram}

Compactifying the Kruskal coordinates via
\begin{equation}
	\bar u = \tan^{-1}\frac{\tilde u}{4M},
	\qquad \bar v = \tan^{-1}\frac{\tilde v}{4M},
\end{equation}
maps the entire spacetime into a finite region. This choice is not
unique; any monotonic functions would preserve the causal relations.
The resulting Penrose diagram represents the full causal structure of
the spacetime in a finite picture, with null rays propagating at
$45°$ throughout.

% [FIGURE: Penrose diagram of the Schwarzschild spacetime showing
%  regions I-IV, the singularity, the horizon, and the asymptotic
%  regions scri^+, scri^-, i^+, i^-, i^0]

The causal structure encoded in the diagram is physically important.
An observer in region I can receive signals from region I and from
region IV, but cannot communicate with the interior (region II) or
with another exterior observer in region III. The singularity at $r=0$
is a spacelike surface in regions II, meaning it lies in the future of
every observer who crosses the horizon, and cannot be avoided.

The Penrose diagram for 2D Minkowski spacetime is obtained by a
similar construction. In null coordinates $u = t-x$, $v = t+x$, the
metric is $ds^2 = -du\,dv$. The compactification
\begin{equation}
	u' = 2\tan^{-1} u, \qquad v' = 2\tan^{-1} v,
	\qquad -\pi \leq u', v' \leq \pi,
\end{equation}
maps the full Minkowski space to a finite diamond. The conformal factor
is
\begin{equation}
	\Omega^2 = \frac{1}{4}\sec^2\!\frac{u'}{2}\sec^2\!\frac{v'}{2},
\end{equation}
and after a conformal rescaling $g_{\mu\nu}\to\Omega^{-2}g_{\mu\nu}$
the metric takes the same form $d\tilde s^2 = du'\,dv'$ over the
compact region. The boundaries of the diagram correspond to future and
past null infinity $\mathscr{I}^\pm$, future and past timelike infinity
$i^\pm$, and spacelike infinity $i^0$. Null rays terminate on
$\mathscr{I}^\pm$ because $u = \tan(u'/2)\to\pm\infty$ as
$u'\to\pm\pi$. In 4D the same construction applies, with each interior
point representing a 2-sphere, except on the symmetry axis and at $i^0$.

% [FIGURE: Penrose diagram for 2D and 4D Minkowski spacetime]
